\documentclass[12pt,a4paper]{article}
\usepackage{fontspec}
\usepackage{geometry}
\usepackage{enumitem}
\usepackage{fancyhdr}
\usepackage{lastpage}
\usepackage{hyperref}

\setmainfont{IBM Plex Sans Devanagari}[Scale=1.0]

\geometry{margin=2.5cm}
\pagestyle{fancy}
\fancyhf{}
\rhead{Page \thepage\ of \pageref{LastPage}}
\lhead{ECO104 - Quiz (hi)}
\renewcommand{\headrulewidth}{0.4pt}

\title{Quiz: ECO104}
\author{Language: hi}
\date{\today}

\begin{document}

\maketitle
\thispagestyle{fancy}

\section*{Instructions}
Answer all questions by selecting the best option (A, B, C, or D). The answer key is provided at the end of this document.

\bigskip


\subsection*{Question 1}
निम्नलिखित में से कौन सा पैसे के तीन आवश्यक कार्यों में से एक *नहीं* है?

\begin{enumerate}[label=\Alph*., leftmargin=*]
\item ऊर्जा का स्रोत
\item हिसाब की इकाई
\item Exchange का माध्यम
\item मूल्य संचय (Store of Value)
\end{enumerate}

\bigskip


\subsection*{Question 2}
निम्नलिखित में से क्या बार्टर प्रणाली में एक सफल लेन-देन के लिए ज़रूरी है?

\begin{enumerate}[label=\Alph*., leftmargin=*]
\item एक आम खाता इकाई
\item न्यूनतम शेष राशि की आवश्यकता
\item एक केंद्रीय बैंकिंग प्राधिकरण
\item दोहरी इच्छाओं का संयोग
\end{enumerate}

\bigskip


\subsection*{Question 3}
पैसे और सोने जैसी चीज़ों के बीच ऐतिहासिक कनेक्शन क्या है?

\begin{enumerate}[label=\Alph*., leftmargin=*]
\item सोना आसानी से बाँटा जा सकता है
\item सोना दुर्लभ है और मूल्य संचय के रूप में व्यापक रूप से स्वीकार किया जाता है।
\item सोना पैसे को और भी आकर्षक बना देता है
\item सोना अनोखा है
\end{enumerate}

\bigskip


\subsection*{Question 4}
धातु के सिक्कों का एक फायदा यह था कि वे वस्तु-मुद्रा (commodity money) के मुकाबले ज़्यादा टिकाऊ और आसानी से ले जाने वाले होते थे।

\begin{enumerate}[label=\Alph*., leftmargin=*]
\item यह ज्यादा आसानी से ले जाने योग्य था
\item यह नकली बनाने के लिए आसान नहीं था
\item इसका आंतरिक मूल्य ज़्यादा था
\item वो और ज़्यादा टिकाऊ था
\end{enumerate}

\bigskip


\subsection*{Question 5}
एक कारण यह है कि फिएट करेंसी से केंद्रीकृत नियंत्रण क्यों होता है?

\begin{enumerate}[label=\Alph*., leftmargin=*]
\item यह विकेंद्रित वित्तीय प्रणालियों को प्रोत्साहित करता है।
\item यह महँगाई के खिलाफ मजबूत है।
\item यह सोने और चांदी जैसी भौतिक वस्तुओं द्वारा समर्थित है।
\item यह सरकारों को पैसे Supply और ब्याज दरों पर नियंत्रण करने की अनुमति देता है।
\end{enumerate}

\bigskip


\subsection*{Question 6}
मुद्रास्फीति (इन्फ्लेशन) समय के साथ वेतनभोगियों की खरीदारी क्षमता को कैसे प्रभावित करती है?

\begin{enumerate}[label=\Alph*., leftmargin=*]
\item महंगाई सिर्फ उन लोगों की खरीदने की ताकत को प्रभावित करती है जो फाइनेंशियल मार्केट्स में निवेश करते हैं।
\item महँगाई से खरीदारी की ताकत बढ़ती है क्योंकि मजदूरी कीमतों से ज्यादा तेजी से बढ़ती है।
\item महँगाई खरीदने की ताकत को कम कर देती है क्योंकि मजदूरी अक्सर असली महँगाई दर से पीछे रह जाती है।
\item महंगाई का खरीदारी की क्षमता पर कोई असर नहीं पड़ता।
\end{enumerate}

\bigskip


\subsection*{Question 7}
आज घर खरीदना 1980s के मुकाबले इतना मुश्किल क्यों है?

\begin{enumerate}[label=\Alph*., leftmargin=*]
\item प्रॉपर्टी की कमी और दबे हुए ब्याज दरों की वजह से कर्ज़ ज़्यादा आकर्षक लग रहा है।
\item जीवनयापन की बढ़ती लागत और जनसंख्या वृद्धि
\item महँगाई और किराया दरें
\item दबाए हुए ब्याज दरें कर्ज़ को और आकर्षक बना देती हैं और मुद्रास्फीति
\end{enumerate}

\bigskip


\subsection*{Question 8}
ज्यादा पैसा छापने का करेंसी की खरीदने की ताकत पर क्या असर पड़ता है?

\begin{enumerate}[label=\Alph*., leftmargin=*]
\item कोई असर नहीं
\item ज़्यादा पैसा = ज़्यादा ख़रीदारी की ताक़त
\item मुद्रा की क्रय शक्ति में वृद्धि
\item मुद्रा की क्रय शक्ति में गिरावट
\end{enumerate}

\bigskip


\subsection*{Question 9}
नए Bitcoin कैसे बनाए जाते हैं और परिचालन में लाए जाते हैं?

\begin{enumerate}[label=\Alph*., leftmargin=*]
\item एक केंद्रीय प्राधिकरण द्वारा जारी करने के माध्यम से
\item उन्हें लेन-देन शुल्क के रूप में प्राप्त करके
\item Bitcoin एक्सचेंज से खरीदकर
\item Blockchain में नए लेन-देन जोड़ने के इनाम के रूप में Mining के माध्यम से
\end{enumerate}

\bigskip


\subsection*{Question 10}
द टाइम्स 03/जनवरी/2009 चांसलर प्रथम जीडब्ल्यू-0 ब्लॉक में बैंकों के लिए दूसरे बेलआउट के कगार पर, इस उद्धरण का क्या महत्व है?

\begin{enumerate}[label=\Alph*., leftmargin=*]
\item यह इस बात का एक बेहतरीन उदाहरण है कि किस तरह बेलआउट से हमारे समाज को लाभ हुआ है
\item यह हमें Bitcoin की Blockchain में संदेशों को संग्रहीत करने की क्षमता दिखाता है
\item यह टाइम्स अख़बार के प्रति Satoshi की सराहना को दर्शाता है
\item यह हमारी वित्तीय प्रणाली की घाटे को सामाजिक बनाने की क्षमता को उजागर करता है
\end{enumerate}

\bigskip


\subsection*{Question 11}
Bitcoin नेटवर्क में नोड्स का क्या उद्देश्य है?

\begin{enumerate}[label=\Alph*., leftmargin=*]
\item दूसरे माइनरों से मुकाबला करने के लिए
\item लेन-देन प्रोसेस करने के लिए
\item Blockchain को बनाए रखने और सुरक्षित करने के लिए
\item Bitcoin कोडबेस में बदलाव सुझाना और लागू करना
\end{enumerate}

\bigskip


\subsection*{Question 12}
Bitcoin और पारंपरिक मुद्राओं में क्या अंतर है जो केंद्रीय बैंकों या सरकारों द्वारा नियंत्रित होती हैं?

\begin{enumerate}[label=\Alph*., leftmargin=*]
\item जीडब्ल्यू-1 को माइनर्स, नोड्स और डेवलपर्स के बीच सहयोग की जरूरत है
\item Bitcoin एक विकेंद्रीकृत सहमति प्रणाली पर निर्भर करता है।
\item उन सभी को
\item जीडब्ल्यू-2 नेटवर्क में किसी को भी भाग लेने और महत्वपूर्ण भूमिकाएं निभाने की अनुमति देता है
\end{enumerate}

\bigskip


\subsection*{Question 13}
Bitcoin के नेटवर्क में पारंपरिक वित्तीय प्रणालियों पर क्या फायदा है?

\begin{enumerate}[label=\Alph*., leftmargin=*]
\item महंगाई दरों में बढ़ोतरी
\item बिचौलियों पर निर्भरता
\item विकेंद्रीकरण और सुरक्षा
\item लेन-देन पर केंद्रीकृत नियंत्रण
\end{enumerate}

\bigskip


\subsection*{Question 14}
Bitcoin कैसे युद्धग्रस्त या अस्थिर देशों में लोगों को फायदा पहुंचाता है?

\begin{enumerate}[label=\Alph*., leftmargin=*]
\item यह एक भौतिक संपत्ति प्रदान करता है जिसका मूल्य बढ़ सकता है
\item यह मुद्रास्फीति के खिलाफ सुरक्षा प्रदान करता है
\item यह वैल्यू को एक डिसेंट्रलाइज़्ड और सिक्योर तरीके से स्टोर करने की सुविधा देता है जिसमें चोरी का ज़्यादा रेज़िस्टेंस होता है।
\item यह पारंपरिक बैंकिंग सेवाओं तक पहुंच की अनुमति देता है
\end{enumerate}

\bigskip


\subsection*{Question 15}
Bitcoin लेनदेन में प्राइवेट कीज़ की क्या भूमिका है?

\begin{enumerate}[label=\Alph*., leftmargin=*]
\item वे लेन-देन के डेटा को एन्क्रिप्ट करने के लिए इस्तेमाल होते हैं।
\item वे लेन-देन की असली होने की पुष्टि करने के लिए इस्तेमाल किए जाते हैं।
\item वे पारदर्शिता के लिए सार्वजनिक रूप से साझा किए जाते हैं।
\item वे लेन-देन पर हस्ताक्षर करने और उन्हें अधिकृत करने के लिए इस्तेमाल किए जाते हैं।
\end{enumerate}

\bigskip


\subsection*{Question 16}
गैर-कस्टोडियल वॉलेट्स का Bitcoin स्टोरेज के लिए मुख्य फायदा क्या है?

\begin{enumerate}[label=\Alph*., leftmargin=*]
\item कम लेन-देन शुल्क
\item पासवर्ड की ज़रूरत नहीं
\item बेहतर सुरक्षा और गोपनीयता
\item Bitcoin तक आसान पहुँच
\end{enumerate}

\bigskip


\subsection*{Question 17}
स्व-संरक्षण लेते समय, नुकसान के जोखिम को न्यूनतम करने के लिए आपको सबसे पहले क्या करना चाहिए?

\begin{enumerate}[label=\Alph*., leftmargin=*]
\item अपने Wallet पुनर्प्राप्ति seed वाक्यांश का बैकअप लें
\item अपनी निजी कुंजियाँ किसी विश्वसनीय पक्ष के साथ साझा करें
\item अपने पुनर्प्राप्ति seed वाक्यांश को नोट करें और इसे अपने कंप्यूटर पर संग्रहीत करें
\item दो-कारक प्रमाणीकरण (2FA) अक्षम करें
\end{enumerate}

\bigskip


\subsection*{Question 18}
पाठ में उल्लिखित स्टेबलकॉइन के मुख्य प्रकार क्या हैं?

\begin{enumerate}[label=\Alph*., leftmargin=*]
\item फिएट-बैक्ड, कमोडिटी-बैक्ड और एल्गोरिदमिक स्टेबलकॉइन्स।
\item केंद्रीकृत और विकेंद्रीकृत स्थिरमुद्राएँ
\item बैंकों द्वारा समर्थित स्टेबलकॉइन और व्यक्तियों द्वारा समर्थित स्टेबलकॉइन
\item डिजिटल और भौतिक स्टेबलकॉइन्स
\end{enumerate}

\bigskip


\subsection*{Question 19}
2015 में ग्रीक सरकार के दिवालिया होने के संकट के दौरान लोगों को क्या एक बड़ी समस्या का सामना करना पड़ा था?

\begin{enumerate}[label=\Alph*., leftmargin=*]
\item डिजिटल बैंकिंग सेवाओं तक सीमित पहुँच
\item भौतिक मुद्रा (नकदी) की कमी
\item जीडब्ल्यू-0 जैसी डिजिटल करेंसी की अस्थिरता
\item सरकार द्वारा थोपा गया दबंग बैंकिंग सुधार।
\end{enumerate}

\bigskip


\subsection*{Question 20}
स्टेबलकॉइन, जैसे कि USDt, अपनी वैल्यू को रिलेटिवली स्टेबल कैसे बनाए रखते हैं?

\begin{enumerate}[label=\Alph*., leftmargin=*]
\item वे दूसरी क्रिप्टोकरेंसी की तरह उतार-चढ़ाव के शिकार होते हैं।
\item वे सुरक्षित लेन-देन के लिए Blockchain टेक्नोलॉजी पर भरोसा करते हैं।
\item वे किसी केंद्रीय प्राधिकरण या सरकार द्वारा समर्थित हैं।
\item वे एक स्थिर संपत्ति या मुद्रा से जुड़े हुए हैं।
\end{enumerate}

\bigskip


\subsection*{Question 21}
टेदर के विभिन्न स्टेबलकॉइन्स अस्थिर अर्थव्यवस्थाओं में व्यक्तियों को वित्तीय स्वतंत्रता कैसे देते हैं?

\begin{enumerate}[label=\Alph*., leftmargin=*]
\item बचत पर ज़्यादा ब्याज देकर
\item भौतिक नकदी की जरूरत को खत्म करके
\item स्थिर Exchange दरों वाली मुद्राओं तक पहुंच प्रदान करके
\item वैश्विक लेन-देन को तुरंत संभव बनाकर
\end{enumerate}

\bigskip


\subsection*{Question 22}
स्टेबलकॉइन रेमिटेंस से जुड़े खर्चों को कैसे कम करते हैं?

\begin{enumerate}[label=\Alph*., leftmargin=*]
\item मुफ्त पैसे ट्रांसफर की सुविधा देकर
\item लोगों को पारंपरिक सेवाओं से बेहतर Exchange दरें देकर
\item फिएट करेंसी को डिजिटाइज़ करके
\item पारंपरिक सेवाओं से कम फीस लेकर
\end{enumerate}

\bigskip


\subsection*{Question 23}
दुनिया भर में कितने देश ऐसे हैं जहाँ महिलाओं को अपने नाम से बैंक खाता खोलने पर रोक है?

\begin{enumerate}[label=\Alph*., leftmargin=*]
\item १००
\item ३६
\item 75
\item 72
\end{enumerate}

\bigskip


\subsection*{Question 24}
Halving, Bitcoin के संदर्भ में क्या है?

\begin{enumerate}[label=\Alph*., leftmargin=*]
\item Bitcoin के वैश्विक अपनाव को बढ़ाने की एक प्रक्रिया
\item एक शब्द जिसका इस्तेमाल उस घटना को बताने के लिए किया जाता है जब Bitcoin की कीमत आसमान छूने लगती है
\item Bitcoin के कोड में एक फीचर जो हर चार साल में Mining का इनाम आधा कर देता है
\item एक ऐसी व्यवस्था जो Bitcoin के मूल्य को गिरने से रोकती है
\end{enumerate}

\bigskip


\subsection*{Question 25}
यह तर्क कि Bitcoin बेकार है क्योंकि यह किसी चीज़ से समर्थित नहीं है, गलत क्यों है?

\begin{enumerate}[label=\Alph*., leftmargin=*]
\item जीडब्ल्यू-1 सोने से बैक्ड है, जिससे इसकी अपनी एक अंतर्निहित वैल्यू होती है।
\item Bitcoin की कीमत इसकी कम उपलब्धता और सुरक्षा फीचर्स से तय होती है।
\item जीडब्ल्यू-3 की कीमत सरकारों और बैंकों द्वारा इसकी स्वीकृति पर आधारित है।
\item वैल्यू सब्जेक्टिव होती है और लोगों की मान्यताओं और अनुभवों पर आधारित होती है।
\end{enumerate}

\bigskip


\subsection*{Question 26}
क्रिप्टोकरेंसी की दुनिया में Blockchain ट्राइलेमा क्या है?

\begin{enumerate}[label=\Alph*., leftmargin=*]
\item Bitcoin की एक खासियत यह है कि यह स्केलेबिलिटी के बजाय विकेंद्रीकरण और सुरक्षा को प्राथमिकता देता है।
\item क्रिप्टोकरेंसी के रास्ते और मूल्यों को समझने में मुश्किल
\item Bitcoin की तकनीक की सीमाएँ अन्य क्रिप्टोकरेंसी की तुलना में
\item क्रिप्टोकरेंसी में स्केलेबिलिटी, सुरक्षा और विकेंद्रीकरण के बीच ट्रेड-ऑफ (समझौता) होता है।
\end{enumerate}

\bigskip


\subsection*{Question 27}
Bitcoin के कुल लेन-देन की मात्रा का कितना प्रतिशत अवैध गतिविधियों के लिए इस्तेमाल किया गया है?

\begin{enumerate}[label=\Alph*., leftmargin=*]
\item २-४%
\item 1% से कम
\item 5% से ज़्यादा
\item यह अज्ञात है
\end{enumerate}

\bigskip


\subsection*{Question 28}
क्या कोई भी Bitcoin कोड को कॉपी कर सकता है?

\begin{enumerate}[label=\Alph*., leftmargin=*]
\item हाँ, लेकिन तुम Bitcoin नेटवर्क का वैल्यू कॉपी नहीं कर सकते।
\item नहीं, क्योंकि Bitcoin कोड क्लोज्ड-सोर्स है
\item नहीं और आप Bitcoin नेटवर्क का मान कॉपी कर सकते हैं
\item हाँ और तुम Bitcoin नेटवर्क का वैल्यू भी कॉपी करोगे
\end{enumerate}

\bigskip


\subsection*{Question 29}
Bitcoin, Mining के जरिए ऊर्जा का उपभोग करने के पीछे एक प्रेरणा क्या है?

\begin{enumerate}[label=\Alph*., leftmargin=*]
\item Bitcoin को इस्तेमाल करने में और महंगा बनाने के लिए
\item नए यूज़र्स को नेटवर्क में शामिल होने से हतोत्साहित करने के लिए
\item केंद्रीकरण के खिलाफ एक ताकतवर सुरक्षा तंत्र प्रदान करने के लिए
\item इसकी वैल्यू कम करने के लिए एनर्जी कॉस्ट बढ़ाना
\end{enumerate}

\bigskip


\subsection*{Question 30}
एक स्टेबलकॉइन पर भरोसा करने से पहले अच्छी तरह रिसर्च करना क्यों ज़रूरी है? एक वजह ये है कि कुछ स्टेबलकॉइन्स के पीछे पर्याप्त संपत्ति या रिज़र्व नहीं होता, जिससे उनकी वैल्यू अचानक गिर सकती है!

\begin{enumerate}[label=\Alph*., leftmargin=*]
\item कुछ स्टेबलकॉइन सरकार द्वारा रेगुलेटेड होते हैं
\item स्टेबलकॉइन हमेशा रिजर्व से पूरी तरह बैक्ड होते हैं
\item कुछ स्टेबलकॉइन में पारदर्शिता और सावधानी की कमी होती है
\item स्टेबलकॉइन में कोई जोखिम नहीं होता है
\end{enumerate}

\bigskip


\subsection*{Question 31}
स्टेबलकॉइन्स में शॉर्ट-टर्म प्राइस वोलेटिलिटी का मुख्य कारण क्या है?

\begin{enumerate}[label=\Alph*., leftmargin=*]
\item एक्सचेंजों पर तरलता का स्तर कम है
\item निकासी की मांग पूरी करने के लिए रिजर्व की कमी
\item स्टेबलकॉइन का डिपेगिंग
\item हैक्स जिसने स्टेबलकॉइन को कमज़ोर कर दिया
\end{enumerate}

\bigskip


\subsection*{Question 32}
एक स्थिर मुद्रा (stablecoin) अपने निर्धारित मूल्य से ऊपर किस परिस्थिति में कारोबार हो सकती है?

\begin{enumerate}[label=\Alph*., leftmargin=*]
\item जब यूजर्स की तरफ से जोरदार सेलिंग प्रेशर होता है
\item जब स्टेबलकॉइन की मांग Exchange की क्षमता से अधिक हो जाती है
\item जब स्टेबलकॉइन जारीकर्ता फंड फ्रीज कर देता है
\item जब स्टेबलकॉइन पूरी तरह से नकद भंडार द्वारा समर्थित होता है
\end{enumerate}

\bigskip


\subsection*{Question 33}
केंद्रीकृत स्टेबलकॉइन का एक फायदा क्या है, भले ही वे फंड्स को फ्रीज़ कर सकते हों?

\begin{enumerate}[label=\Alph*., leftmargin=*]
\item वे पूरी तरह से विकेंद्रीकरण और सुरक्षा प्रदान करते हैं
\item वे कानून प्रवर्तन से स्वतंत्र रूप से काम करते हैं
\item वे कुछ खास स्थितियों में यूजर्स की सुरक्षा के लिए तुरंत जवाब दे सकते हैं।
\item वे किसी भी नियम या हस्तक्षेप के अधीन नहीं हैं
\end{enumerate}

\bigskip


\newpage
\section*{Answer Key}

\begin{enumerate}
\item \textbf{Question 1:} A
\\ \textit{पैसे के तीन मुख्य काम हैं - मूल्य का संचय, विनिमय का माध्यम (Exchange),
और लेखा की इकाई। ऊर्जा का स्रोत पैसे का काम नहीं है।}

\item \textbf{Question 2:} D
\\ \textit{बार्टर सिस्टम में, दो पार्टियाँ सीधे सामान या सेवाओं का आदान-प्रदान करके लेन-देन करती हैं।
पैसे के बिना सेवाएँ। एक सफल लेन-देन के लिए दोहरे संयोग की आवश्यकता होती है।
चाहिए, मतलब दोनों पक्षों के पास कुछ ऐसा होना चाहिए जो दूसरा पार्टी ट्रेड करना चाहता हो।}

\item \textbf{Question 3:} B
\\ \textit{सोने की प्राकृतिक रूप से दुनिया भर में कमी और इसमें किए जाने वाले बड़े निवेश को देखते हुए
मेरे लिए, सोना खोदना, साफ करना और स्टोर करना हमेशा से प्रीमियम दर पर ही ट्रेड हुआ है...
पैसे के अन्य रूप।}

\item \textbf{Question 4:} D
\\ \textit{धातु के पैसे, खासकर सोना और चाँदी, सामान के पैसे से ज़्यादा टिकाऊ होते थे।
कीमती धातुएं समय के साथ खराब होने से कम प्रभावित होती हैं, जिससे वे आदर्श होती हैं...
मूल्य का लंबे समय तक भंडारण। यह टिकाऊपन भी उनकी उपयोगिता को बढ़ाता था।
Exchange का एक माध्यम, लेखांकन की एक इकाई, और मूल्य का भंडार।}

\item \textbf{Question 5:} D
\\ \textit{फिएट करेंसी सरकारों को पैसे Supply को कंट्रोल करने की ताक़त देती है और
ब्याज दरें, जिससे वे आर्थिक गतिविधि को प्रभावित कर सकते हैं और इसका परिणाम होता है
केंद्रीकृत नियंत्रण में।}

\item \textbf{Question 6:} C
\\ \textit{महंगाई से उपभोक्ता कीमतों में लगातार बढ़ोतरी होती है, जिससे खरीदने की क्षमता कम हो जाती है।
मजदूरों की ताकत कम हो जाती है क्योंकि उनकी मजदूरी महंगाई की असली दर से पीछे रह जाती है।
इससे कमाई करने वालों के लिए अपने लंबे समय के बचत के लक्ष्यों को हासिल करना मुश्किल हो जाता है
और अपने जीवन स्तर को बनाए रखें।}

\item \textbf{Question 7:} D
\\ \textit{जब मुद्रास्फीति हमारी मुद्रा की क्रय शक्ति को कम कर रही है, तो यह
सेविंग करना मुश्किल हो रहा है, दबे हुए ब्याज़ दर कर्ज़े की खपत बढ़ा रहे हैं
ज़्यादा आकर्षक। जब पूंजी की लागत इतनी सस्ती होती है, तो लोग ज़रूरत से ज़्यादा उधार ले लेते हैं।
उनके पास जो भी पैसा है, वो संपत्तियों में डाल रहे हैं, जिससे कीमतें बढ़ रही हैं।}

\item \textbf{Question 8:} D
\\ \textit{जब केंद्रीय बैंक पैसा छापने का फैसला करते हैं, तो वे बस ऐसे ही ज्यादा वैल्यू नहीं बना सकते।
इस नए छपे हुए पैसे का कुछ मूल्य होने के लिए, इसकी कीमत कहीं से तो आनी चाहिए...
पिछले मुद्रा धारकों।
Supply के पैसे को एक पिज़्ज़ा की तरह समझो, और कल्पना करो कि इसे चार टुकड़ों में काटा गया है।
Supply का पैसा दोगुना करना, पिज़्ज़ा की मात्रा दोगुना करने के बराबर नहीं होगा।
इस्टेड, ये उन चार स्लाइस को आधा काटकर बनाने के बराबर होगा
आठ टुकड़े। हमें कोई अतिरिक्त पिज़्ज़ा नहीं मिला है। हमारे पास बस ज़्यादा टुकड़े हैं,
हर नोट छोटे आकार का होता है। इसलिए, जब हम ज्यादा पैसा छापते हैं, तो उसकी कीमत कम हो जाती है।
वो पहले से ही मौजूद है।}

\item \textbf{Question 9:} D
\\ \textit{नए Bitcoin को Mining नामक प्रक्रिया के जरिए चलन में लाया जाता है।
माइनर्स एक-दूसरे के साथ Blockchain में नए ट्रांजैक्शन जोड़ने के लिए प्रतिस्पर्धा करते हैं।
Exchange में अपने काम के लिए, माइनर्स को नए बनाए गए Bitcoin से इनाम मिलता है।
यह प्रक्रिया Bitcoin नेटवर्क पर लेन-देन को सत्यापित और सुरक्षित भी करती है।}

\item \textbf{Question 10:} D
\\ \textit{यह टाइम्स अख़बार के एक लेख का संदर्भ है [लेख](https://www.thetimes.co.uk/article/chancellor-alistair-darling-on-brink-of-second-bailout-for-banks-n9l382mn62h) जो Satoshi की चिंताओं को दर्शाता है कि बैंक जोखिम भरा व्यवहार कर रहे हैं, जिसका उन पर कोई खास असर नहीं होगा और यह घाटा मुद्रा धारकों के बीच साझा किया जाएगा।}

\item \textbf{Question 11:} C
\\ \textit{नोड्स Bitcoin नेटवर्क की सुरक्षा और मजबूती के लिए बेहद जरूरी हैं।
वे सिस्टम के गेटकीपर की तरह काम करते हैं, यह सुनिश्चित करते हुए कि लेन-देन
सही तरीके से और नियमों के अनुसार प्रोसेस किए जाते हैं।}

\item \textbf{Question 12:} C
\\ \textit{Bitcoin और पारंपरिक फिएट करेंसी में फर्क यह है कि यह एक...}

\item \textbf{Question 13:} C
\\ \textit{Bitcoin का विकेंद्रीकृत नेटवर्क एक सुरक्षित और कुशल तरीका प्रदान करता है
बीच में किसी बिचौलिए या तीसरे पक्ष की जरूरत के बिना लेन-देन करें। यह
सेंसरशिप और लेन-देन में अनधिकृत बदलाव का खतरा कम करता है,
और साथ ही प्राइवेसी भी बढ़ाता है।}

\item \textbf{Question 14:} C
\\ \textit{युद्धग्रस्त या अस्थिर देशों में, लोगों को मुश्किलों का सामना करना पड़ सकता है...
पारंपरिक बैंकिंग सेवाओं का उपयोग करो या चोरी के खतरे का सामना करो
भौतिक संपत्ति की ज़ब्ती। Bitcoin व्यक्तियों को मूल्य संग्रहित करने की अनुमति देता है।
एक विकेंद्रित और सुरक्षित तरीके से, बिना किसी बिचौलिए की जरूरत के
या फिर भौतिक संपत्ति जो महंगाई, चोरी या ज़ब्त जैसे जोखिमों के प्रति संवेदनशील हो।}

\item \textbf{Question 15:} D
\\ \textit{प्राइवेट कीज़ Bitcoin लेन-देन का एक जरूरी हिस्सा हैं।
वे लेन-देन पर हस्ताक्षर करने और उन्हें अधिकृत करने के लिए इस्तेमाल किए जाते हैं, जो सबूत प्रदान करते हैं
जीडब्ल्यू-2 की और यह सुनिश्चित करना कि केवल मालिक ही अपने जीडब्ल्यू-3 खर्च कर सकता है।}

\item \textbf{Question 16:} C
\\ \textit{नॉन-कस्टोडियल वॉलेट्स वो होते हैं जहां आप अपने फंड्स के एकमात्र मालिक होते हैं,
इसका मतलब है कि आपके पास अपनी प्राइवेट कीज़ पर पूरा कंट्रोल है। इसलिए,
नॉन-कस्टोडियल वॉलेट्स ज़्यादा सुरक्षा और प्राइवेसी देते हैं
कस्टोडियल वॉलेट पर क्योंकि आप ही अपने पैसों की सुरक्षा के लिए ज़िम्मेदार होते हैं।}

\item \textbf{Question 17:} A
\\ \textit{सबसे पहले और सबसे महत्वपूर्ण बात, स्व-संरक्षण लेते समय अपने Wallet का बैकअप लें। हार्डवेयर वॉलेट में रिकवरी seed वाक्यांश होता है, शब्दों का एक सेट जिसका उपयोग आपके डिवाइस के खो जाने या क्षतिग्रस्त होने की स्थिति में आपकी निजी कुंजियों को पुनर्प्राप्त करने के लिए किया जा सकता है। इस seed वाक्यांश की एक भौतिक प्रतिलिपि बनाएँ, जैसे कि धातु [seed प्लेट](https: //coincodex.com/article/23147/best-metal-crypto-wallets-for-seed-phrase-storage/), और इसे सुरक्षित स्थान पर संग्रहीत करें। इस रिकवरी seed वाक्यांश को सुरक्षित रखना महत्वपूर्ण है। इसे कभी भी किसी के साथ साझा न करें।}

\item \textbf{Question 18:} A
\\ \textit{पिछले अध्यायों में हमने स्टेबलकॉइन के तीन मुख्य प्रकार देखे हैं;
फिएट-बैक्ड स्टेबलकॉइन, जो पारंपरिक मुद्राओं द्वारा समर्थित होते हैं
अमेरिकी डॉलर की तरह; कमोडिटी-बैक्ड स्टेबलकॉइन्स, जो समर्थित होते हैं
सोना या तेल जैसी वस्तुएं; और एल्गोरिदमिक स्टेबलकॉइन्स, जो निर्भर करते हैं...
एक नियमों या एल्गोरिदम के सेट पर उनके मूल्य को बनाए रखने के लिए। हालांकि,
यह बताता है कि सफल एल्गोरिदमिक स्टेबलकॉइन अभी तक बाजार में नहीं आए हैं।}

\item \textbf{Question 19:} D
\\ \textit{2015 में यूनानी सरकार के दिवालिया होने के संकट के दौरान, बैंक बंद हो गए थे,
ATM से पैसे निकालने पर पाबंदी थी, और लोगों का एक बड़ा हिस्सा...
सरकार ने बैंक जमाओं से पैसे निकाल लिए ताकि वो अपने फिजूलखर्ची वाले कामों को पूरा कर सके।
इस घटना ने पारंपरिक बैंकिंग प्रणाली में पैसे रखने से जुड़े जोखिमों को उजागर किया।}

\item \textbf{Question 20:} D
\\ \textit{स्टेबलकॉइन्स अपनी वैल्यू को स्थिर रखने के लिए किसी चीज़ से जुड़े (पेग्ड) होते हैं।
एक और स्थिर संपत्ति या मुद्रा में, जैसे कि अमेरिकी डॉलर या सोना।
यह पेगिंग दूसरे मुद्राओं के साथ जुड़ी उतार-चढ़ाव को कम करने में मदद करती है।
डिजिटल एसेट्स जैसे Bitcoin.}

\item \textbf{Question 21:} C
\\ \textit{टेदर के विभिन्न स्टेबलकॉइन अमेरिकी डॉलर जैसी स्थिर संपत्तियों से जुड़े हुए हैं।
डॉलर या सोना। यह स्थिरता अस्थिर अर्थव्यवस्थाओं में लोगों को
अपने पैसे को एक ऐसी मुद्रा में ले जाना जो उसी के अधीन नहीं है
महंगाई दरों को नियंत्रित करके, उनकी बचत की कीमत बनाए रखने में मदद मिलती है।}

\item \textbf{Question 22:} D
\\ \textit{स्टेबलकॉइन, पारंपरिक मनी ट्रांसफर सर्विसेज के विपरीत, काफी
कम फीस क्योंकि इन्हें बैंक या पैसे जैसे बिचौलियों की ज़रूरत नहीं होती
ट्रांसफर सर्विसेज।}

\item \textbf{Question 23:} D
\\ \textit{दुनिया के 72 देशों में, महिलाओं को खाता खोलने की अनुमति नहीं है
अपने नाम से बैंक खाते।}

\item \textbf{Question 24:} C
\\ \textit{Halving, Bitcoin के कोड में एक फीचर है जो Mining को काट देता है
हर चार साल में लगभग आधा इनाम मिलता है। यह घटना प्रभावित करती है
खनिकों की आमदनी बढ़ती है और Bitcoin की कीमत पर ऊपर की ओर दबाव पड़ता है।}

\item \textbf{Question 25:} D
\\ \textit{वैल्यू सब्जेक्टिव होती है और देखने वाले की नज़र में होती है। Bitcoin वैल्यूएबल है।
क्योंकि लोग ऐसा मानते हैं, बिल्कुल वैसे ही जैसे कला कीमती हो सकती है, भले ही...
कैनवास और पेंट से बना होना। यह तर्क कि मुद्रा को किसी चीज़ से समर्थित होना चाहिए...
कुछ मूल्यवान चीज़ ग़लत है, क्योंकि अमेरिकी डॉलर जैसी फ़िएट मुद्राओं का मूल्य
डॉलर भरोसे और विश्वास पर भी आधारित है।}

\item \textbf{Question 26:} D
\\ \textit{Blockchain ट्राइलेम्मा का कॉन्सेप्ट यह है कि क्रिप्टोकरेंसी...
तीन महत्वपूर्ण विशेषताओं में से दो को प्राथमिकता देनी होगी; स्केलेबिलिटी,
सुरक्षा, और विकेंद्रीकरण। Bitcoin विकेंद्रीकरण और सुरक्षा को प्राथमिकता देता है,
जो इसके स्केलेबिलिटी को प्रभावित करता है, लेकिन फिर भी यह वैकल्पिक तरीकों से इसे हासिल कर लेता है।}

\item \textbf{Question 27:} B
\\ \textit{Chainalysis की 2020 की एक रिपोर्ट के मुताबिक, अवैध गतिविधियाँ सिर्फ़
जीडब्ल्यू-1 के सभी लेन-देन का लगभग 0.34% हिस्सा है।}

\item \textbf{Question 28:} A
\\ \textit{जबकि Bitcoin का कोड ओपन सोर्स है और इसे कॉपी किया जा सकता है,
Bitcoin की वैल्यू इसके कोड से नहीं मिलती। यह मिलती है
इसका उपयोगकर्ताओं का नेटवर्क और समय के साथ बना विश्वास व अपनापन।}

\item \textbf{Question 29:} C
\\ \textit{Bitcoin द्वारा Mining के माध्यम से ऊर्जा खपत करने के कारणों में से एक है
केंद्रीकरण के खिलाफ एक शक्तिशाली सुरक्षा तंत्र प्रदान करता है
अनुचित हेरफेर। माइनर्स का वितरित नेटवर्क रोकने में मदद करता है
केंद्रीकृत संस्थाएँ नेटवर्क को अपने कब्जे में लेने या नियमों को अपने फायदे के लिए बदलने से रोकना।}

\item \textbf{Question 30:} C
\\ \textit{सभी स्टेबलकॉइन एक जैसी पारदर्शिता नहीं देते।
इसलिए, अच्छी तरह से रिसर्च करना ज़रूरी है
किसी भी खास स्टेबलकॉइन पर बड़ी रकम सौंपने से पहले।}

\item \textbf{Question 31:} A
\\ \textit{स्टेबलकॉइन में अल्पकालिक कीमतों में उतार-चढ़ाव अक्सर होता है क्योंकि
लिक्विडिटी की समस्याओं के कारण, छोटे एक्सचेंजों में क्षमता की कमी हो सकती है
ग्राहकों की मांग को पूरा करने के लिए, जिससे अस्थायी कीमतों में उतार-चढ़ाव होता है।}

\item \textbf{Question 32:} B
\\ \textit{स्टेबलकॉइन अपने निर्धारित मूल्य से ऊपर कारोबार कर सकते हैं जब
स्टेबलकॉइन की मांग Exchange की क्षमता से ज़्यादा हो गई है
ग्राहकों की मांगों को पूरा करना, आमतौर पर छोटे एक्सचेंजों पर।}

\item \textbf{Question 33:} C
\\ \textit{केंद्रीकृत स्टेबलकॉइन में हस्तक्षेप करने की क्षमता होती है
जब यूजर्स के सर्वोत्तम हित में हो और समग्र स्थिरता के लिए अच्छा हो
सिस्टम की, जैसा कि उन उदाहरणों से पता चलता है जहां जमा हुए फंड थे
उपभोक्ताओं की सुरक्षा के लिए और हैकर्स को फायदा उठाने से रोकने के लिए इस्तेमाल किया जाता है।}


\end{enumerate}

\end{document}
